\section{Conclusões}
\label{sec:conclusoes}

Neste trabalho, analisamos diferentes algoritmos para o Problema do Caixeiro Viajante Euclidiano bidimensional (\texttt{EUC\_2D}), utilizando 15 instâncias da TSPLIB ($N \in [10, 70]$). Ao longo dos experimentos, foram realizadas otimizações graduais com base nas hipóteses levantadas e os resultados foram avaliados por meio de análises estatísticas de desempenho. A seguir, apresentamos as principais conclusões obtidas nas Baterias A--D.

\subsection{Algoritmos Aproximados: \textit{Twice-Around-The-Tree} e Christofides}
\label{sec:conc_aproximados}

Os resultados da Bateria~A (Seção~\ref{sec:res_aproximados}) mostram como a escolha das estruturas de dados pode influenciar o desempenho dos algoritmos polinomiais, além de permitir uma comparação prática entre os dois algoritmos aproximados:

\begin{itemize}
    \item \textbf{\textit{Twice-Around-The-Tree} (\texttt{v0} $\rightarrow$ \texttt{v1}):} a substituição do grafo completo dinâmico, do algoritmo de Kruskal $O(N^2 \log N)$ e do multigrafo duplicado $2T$ pelo Prim denso $\Theta(N^2)$ sobre uma matriz contígua e pelo percurso iterativo em pré-ordem $\Theta(N)$ reduziu consideravelmente o consumo de memória e o tempo de execução. Na instância \texttt{st70} ($N = 70$), por exemplo, o pico de memória passou de $655\text{ KiB}$ para $4{,}1\text{ KiB}$ ($\approx 160\times$), enquanto o tempo mediano caiu de $12{,}2\text{ ms}$ para $0{,}29\text{ ms}$ ($\approx 42\times$). Os resultados de memória ficaram próximos do comportamento teórico esperado, passando de $\Theta(N^2)$ na \texttt{v0} ($R^2 = 0{,}997$) para $\Theta(N)$ na \texttt{v1} ($R^2 = 0{,}996$).

    Em relação à qualidade das soluções, todas as 15 instâncias respeitaram a garantia de $2$-aproximação, com erro médio de $26{,}2\%$ e maior erro de $47{,}1\%$ em \texttt{eil51\_n38}. Mesmo assim, o algoritmo obteve soluções piores que as de Christofides em todas as instâncias e que as do Vizinho Mais Próximo em 12 das 15. Os erros foram maiores nas instâncias com alta entropia espacial ($\rho = 0{,}67$). Por outro lado, não foi encontrada uma relação relevante com o grau máximo da AGM ($\rho = 0{,}03$), contrariando essa expectativa da hipótese $H_{TA}$.

    \item \textbf{Christofides (\texttt{v0} $\rightarrow$ \texttt{v1}):} a remoção do grafo completo também trouxe melhorias para esse algoritmo. Na instância \texttt{st70}, o pico de memória caiu de $655\text{ KiB}$ para $341\text{ KiB}$ e o tempo total passou de $35{,}0\text{ ms}$ para $16{,}5\text{ ms}$. Como previsto na hipótese $H_C$, os ganhos foram limitados pelo emparelhamento perfeito mínimo (Blossom), responsável por $91{,}5\%$ do tempo de execução e por todo o pico de memória da \texttt{v1}.

    Os experimentos também mostraram que o custo do algoritmo está mais relacionado à quantidade de vértices de grau ímpar $|I|$ da AGM do que ao número total de vértices $N$. O ajuste cúbico em $|I|^3$ obteve $R^2 = 0{,}875$, contra $0{,}663$ para $N^3$, enquanto o ajuste de memória em $|I|^2$ atingiu $R^2 = 0{,}999$. Quanto à qualidade das soluções, Christofides apresentou erro médio de $7{,}4\%$ e máximo de $15{,}5\%$, bem abaixo dos $50\%$ previstos pela garantia teórica. Além disso, os resultados se mantiveram relativamente estáveis diante das diferentes características geométricas e estruturais das instâncias ($|\rho| < 0{,}47$).

    \item \textbf{Influência do tamanho das instâncias nas regressões de tempo:} para $N \le 70$, algumas subrotinas otimizadas, como Prim denso, pré-ordem, Hierholzer e Varredura de Graham, executam em menos de $1\text{ ms}$. Nesses casos, o custo fixo do interpretador e as variações causadas pelo escalonador dificultam a observação do crescimento assintótico, o que ajuda a explicar os baixos valores de $R^2$ nas regressões de tempo. Já os resultados de memória ficaram muito próximos das complexidades teóricas ($R^2 \ge 0{,}996$).
\end{itemize}

\subsection{Estratégias de Inicialização (\textit{Warm-Start})}
\label{sec:conc_warmstarts}

Na Bateria~C (Seção~\ref{sec:res_warmstart}), verificamos que as cinco estratégias de inicialização tiveram um custo de execução muito baixo em comparação com a busca exata. Christofides levou, no máximo, $29\text{ ms}$, enquanto a Inserção em Fecho Convexo levou até $4\text{ ms}$. Já a busca exata chegou a durar 30 minutos. Dessa forma, mais importante do que o tempo gasto na inicialização foi a qualidade da solução encontrada, ou seja, o quanto o limite superior inicial $UB_0$ estava próximo do ótimo antes de iniciar a DFS.

\begin{itemize}
    \item \textbf{Inserção em Fecho Convexo}: a otimização com validação preguiçosa reduziu a complexidade para $O(N^2 \log N)$, com um ajuste temporal de $R^2 = 0{,}714$. Além de apresentar um baixo tempo de execução, essa estratégia encontrou as melhores soluções iniciais em 14 das 15 instâncias, com erro médio de apenas $2{,}2\%$, alcançando o ótimo em quatro delas. Também observamos que, quanto maior a proporção de pontos no fecho convexo, menor tende a ser o erro ($\rho = -0{,}70$). Quando utilizada como \textit{warm-start}, a estratégia reduziu em média $41{,}5\%$ dos nós expandidos nas instâncias resolvidas. Em \texttt{eil51\_n25}, por exemplo, o número de nós expandidos passou de 206 milhões para 14 milhões, enquanto o tempo de execução caiu de $341\text{ s}$ para $26{,}1\text{ s}$.

    \item \textbf{Christofides}: entre os algoritmos com garantia formal de aproximação analisados, foi o que encontrou as melhores soluções iniciais, respeitando a garantia de $1{,}5 \cdot OPT$. Os erros variaram entre $1{,}2\%$ e $15{,}5\%$, com média de $7{,}4\%$. Na instância \texttt{ulysses16}, chegou a superar a Inserção em Fecho Convexo, com erro de $1{,}2\%$ contra $1{,}5\%$. Como estratégia de inicialização, reduziu em média $3{,}4\%$ dos nós expandidos e, em \texttt{eil51\_n25}, diminuiu o tempo de execução de $341\text{ s}$ para $295\text{ s}$. Essa redução foi menor porque a DFS geralmente encontra boas soluções logo no início da busca, superando as rotas fornecidas por Christofides. Por outro lado, o algoritmo se mostrou útil nas instâncias que atingiram o limite de 30 minutos (\texttt{TL}). Em comparação com a inicialização fria, os erros finais passaram de $4{,}5\%$ para $2{,}9\%$ em \texttt{ulysses22}, de $9{,}1\%$ para $6{,}8\%$ em \texttt{bays29} e de $18{,}6\%$ para $9{,}1\%$ em \texttt{st70}.

    \item \textbf{Vizinho Mais Próximo}: apesar do baixo tempo de execução e da complexidade $\Theta(N^2)$, a escolha gulosa dos vértices resultou em soluções com erro médio de $16{,}4\%$. Os erros variaram de $5{,}1\%$ em \texttt{eil51\_n38} a $41{,}6\%$ em \texttt{ulysses16}, com tendência a resultados piores em instâncias mais alongadas ($\rho = 0{,}48$ para a razão de aspecto). Na Bateria~C, a redução média dos nós expandidos foi de apenas $0{,}1\%$, pois a estratégia \textit{Best-Child-First} da \texttt{v3} geralmente consegue encontrar soluções melhores logo no início da DFS. Ainda assim, em algumas instâncias, a inicialização trouxe benefícios. Em \texttt{eil51\_n38}, por exemplo, o erro ao atingir o limite de tempo foi de $2{,}7\%$, contra $7{,}2\%$ da inicialização fria.

    \item \textbf{\textit{Twice-Around-The-Tree}}: mesmo com tempo de execução inferior a um milissegundo na \texttt{v1} e garantia teórica de $2 \cdot OPT$, o algoritmo apresentou erros entre $7{,}8\%$ e $47{,}1\%$, com média de $26{,}2\%$. Na prática, as soluções iniciais foram piores do que aquelas encontradas pela DFS logo no início da busca. Por isso, a redução média dos nós expandidos foi de $0{,}0\%$, e a estratégia não trouxe melhorias nas instâncias que atingiram o limite de tempo. Os erros finais foram os mesmos da inicialização fria, como os $18{,}6\%$ em \texttt{st70} e os $9{,}1\%$ em \texttt{bays29}.

    \item \textbf{Inicialização Fria (\textit{Cold Start})}: foi utilizada como base de comparação para as demais estratégias, com complexidade $\Theta(N)$ e sem considerar a distribuição espacial dos vértices. Graças à ordenação dos filhos na \texttt{v3}, seus tempos de execução nas instâncias menores foram próximos aos do Vizinho Mais Próximo e do \textit{Twice-Around-The-Tree}. No entanto, essa abordagem não fornece um limite inicial que favoreça as podas e pode resultar em soluções de menor qualidade quando a execução é interrompida pelo limite de tempo.
\end{itemize}


\subsection{Busca Exata (\textit{Branch-and-Bound} \texttt{v0}--\texttt{v3}) e Aceleração em GPU}
\label{sec:conc_bnb_gpu}

As otimizações realizadas na Bateria~B (Seção~\ref{sec:res_bnb}) e a comparação entre CPU e GPU na Bateria~D (Seção~\ref{sec:res_gpu}) permitiram identificar os principais problemas de desempenho da busca exata implementada em Python.

\begin{itemize}
    \item \textbf{\texttt{v0\_naive} versus \texttt{v1\_gpu} (velocidade e consumo de memória):} na \texttt{v0}, o recálculo de complexidade $O(N^2)$ para cada candidato limitou o processamento a aproximadamente $1\,000\text{ nós/s}$. Com isso, a versão atingiu o limite de tempo de 120~s antes de esgotar a memória disponível, chegando a um pico de $948\text{ MB}$. A vetorização em lotes da \texttt{v1} aumentou a taxa de processamento em $23\times$, chegando a aproximadamente $24\,000\text{ nós/s}$. Em \texttt{burma14}, por exemplo, o tempo caiu de $64{,}3\text{ s}$ para $4{,}37\text{ s}$.

    Apesar desse ganho, o processamento mais rápido fez com que a fila da busca \textit{Best-First} crescesse exponencialmente, chegando à ordem de $10^7$ estados. Como resultado, o limite de 4~GiB de memória foi atingido (\texttt{OOM}) em todas as 11 instâncias maiores, entre $28\text{ s}$ e $69\text{ s}$ de execução.

    \item \textbf{\texttt{v2\_dfs\_memory} versus \texttt{v3\_fast\_bound} (memória e taxa de processamento):} a utilização de DFS \textit{in-place} na \texttt{v2} resolveu o problema de estouro de memória, mantendo a pilha com no máximo $N$ elementos e o consumo do processo estável em $196\text{ MB}$. Entretanto, as alocações de vetores a cada nó ainda limitavam o processamento a cerca de $50\,000\text{ nós/s}$, e a DFS expandia mais estados que a busca \textit{Best-First}.

    A \texttt{v3} reduziu o custo de processamento de cada nó ao combinar o cálculo incremental escalar do limite em $\Theta(1)$, o uso de \textit{bitmask} inteiro e a quebra de simetria. Essas mudanças aumentaram a taxa de processamento em $8\times$ em relação à \texttt{v2}, chegando a aproximadamente $420\,000\text{ nós/s}$. Embora a quebra de simetria possa atrasar a descoberta de boas soluções em algumas instâncias, como \texttt{eil51\_n18}, a \texttt{v3} foi a versão mais rápida em todos os casos resolvidos. Ela também foi a única a resolver \texttt{berlin52\_n20} dentro de 120~s e conseguiu provar a otimalidade em 7 das 15 instâncias dentro do limite de 30 minutos.

    \item \textbf{Uso da GPU:} a opção \texttt{--use-gpu} não trouxe ganhos de velocidade para nenhuma das versões testadas na Bateria~D. A razão mediana entre as taxas de processamento da GPU e da CPU ficou entre $0{,}89$ e $0{,}99$. Além disso, o custo de inicialização do CUDA prejudicou o desempenho nas instâncias menores, como \texttt{eil51\_n10}, que levou $2{,}1\text{ s}$ com GPU e apenas $4\text{ ms}$ com CPU. Para as instâncias analisadas ($N \le 70$), os lotes de processamento não foram grandes o suficiente para aproveitar toda a capacidade da GPU, e o principal gargalo continuou sendo a manipulação da árvore de busca e da fila na CPU. Esse resultado também reforça a decisão de manter o \textit{Twice-Around-The-Tree} executando exclusivamente na CPU.
\end{itemize}

\subsection{Sensibilidade Topológica e Recomendações Finais}
\label{sec:conc_sintese}

A análise de sensibilidade (Seção~\ref{sec:res_sensibilidade}) mostrou que a distribuição espacial dos vértices pode influenciar o desempenho da busca exata mais do que a quantidade de vértices da instância. Distribuições com aglomerados e pontos isolados, caracterizadas por um alto coeficiente de variação das distâncias ao vizinho mais próximo e das arestas da AGM ($\rho = 0{,}90$), além de baixa entropia espacial ($\rho = -0{,}81$), fazem com que o limite inferior de grau~2 na raiz fique mais distante do ótimo. Em \texttt{ulysses22}, por exemplo, a razão $LB_0/OPT$ foi de $0{,}639$, enquanto em \texttt{eil51\_n38} foi de $0{,}936$.

Essa diferença ajuda a explicar por que algumas instâncias menores exigiram mais esforço computacional que outras maiores. A instância \texttt{ulysses16} ($N = 16$) exigiu $340\times$ mais nós e levou $325\text{ s}$, enquanto \texttt{eil51\_n18} ($N = 18$) levou apenas $1{,}6\text{ s}$. De forma semelhante, \texttt{ulysses22} ($N = 22$) ultrapassou o limite de 30 minutos, enquanto \texttt{eil51\_n25} ($N = 25$) foi resolvida em $26{,}1\text{ s}$ utilizando a Inserção em Fecho Convexo.

Com base nos resultados, para instâncias euclidianas pequenas ou com distribuição espacial mais regular ($N \le 25$), a versão \texttt{v3} do \textit{Branch-and-Bound}, inicializada com a Inserção em Fecho Convexo, mostrou-se uma boa escolha por combinar a possibilidade de provar a otimalidade, consumo de memória $\Theta(N)$ e melhor tempo total. Para instâncias maiores ou com muitos aglomerados, a Inserção em Fecho Convexo oferece soluções de boa qualidade em poucos milissegundos, enquanto Christofides é uma alternativa indicada quando é necessário garantir uma aproximação de pior caso de $1{,}5\cdot OPT$.

